\section{Introduction}

The restoration of audio signals corrupted by noise remains a critical challenge in signal processing, particularly when clean reference ground truth is unavailable. Traditional approaches often rely on supervised learning with paired datasets or explicit priors that may not generalize to unseen acoustic environments \cite{lehtinen2018n2n}. While deep generative models have achieved remarkable success in waveform synthesis and denoising tasks by operating directly in the time domain, they frequently struggle to capture fine-grained temporal structures inherent in complex audio signals without leveraging explicit signal processing knowledge. This limitation is particularly evident when addressing artifacts such as wavetable wrap crackle or prolonged residuals that persist near transient events \cite{engel2020ddsp}. To address these challenges, we propose a framework integrating differentiable digital signal processing with automated model design techniques to achieve unsupervised restoration of audio signals while maintaining high fidelity.

Existing methods for hyperparameter optimization and architecture search often treat neural network components as black boxes, neglecting the opportunity to incorporate domain-specific constraints directly into the learning process \cite{elsken2019nas}. Population Based Training (PBT) has demonstrated effectiveness in jointly optimizing model architectures and hyperparameters by maintaining a population of models that compete for computational resources during training \cite{jaderberg2017pbt}. Similarly, differentiable architecture search methods have enabled the automated discovery of neural network structures tailored to specific tasks through continuous relaxation techniques \cite{liu2019darts}. However, these approaches typically lack mechanisms to explicitly encode signal processing priors or handle cases where clean training data is unavailable.

In this work, we make three primary contributions: (1) We introduce a novel metric for evaluating denoising performance that accounts for both residual energy and perceptual quality under unsupervised conditions; (2) We develop an automated search framework combining reinforcement learning with neural architecture search to discover optimal configurations for audio restoration tasks without requiring clean reference signals; (3) We demonstrate the effectiveness of our approach on challenging audio datasets exhibiting complex noise patterns including wavetable wrap artifacts. Our experimental evaluation shows that methods incorporating explicit signal processing knowledge combined with automated design techniques outperform baseline approaches by significant margins in terms of both objective metrics and perceptual quality assessments.

The remainder of this paper is organized as follows: Section~\ref{sec:related_work} reviews related work on audio denoising, neural architecture search, and differentiable signal processing methods. Section~\ref{sec:methodology} details our proposed framework including the metric formulation and automated design procedures. Section~\ref{sec:experiments} presents experimental results comparing our approach against state-of-the-art baselines across various challenging audio scenarios. Finally, Section~\ref{sec:conclusion} summarizes key findings and discusses potential directions for future research in unsupervised audio restoration.
